\section{Introducción}

La resonancia eléctrica es un estado crítico en circuitos de corriente alterna que se presenta cuando la reactancia inductiva y la reactancia capacitiva se igualan en magnitud, lo que resulta en la anulación de la componente reactiva total del sistema. Este fenómeno ocurre a una frecuencia específica, denominada frecuencia de resonancia ($\omega_0$), punto en el cual la impedancia del circuito se vuelve puramente resistiva y alcanza su valor mínimo en configuración serie, permitiendo la máxima transferencia de corriente. \newline

El propósito de este reporte es analizar el material proporcionado por el profesor y validar analíticamente los resultados obtenidos durante la práctica de laboratorio, empleando las técnicas de análisis de circuitos para determinar la selectividad, el factor de calidad y la respuesta en potencia del sistema bajo condiciones de resonancia.